\documentclass[../main.tex]{subfiles}
\ifSubfilesClassLoaded{\setcounter{chapter}{9}}{}
\begin{document}

\chapter{Ensemble Ringkas: Random Forest}

\begin{subcpmk}
  \item \textbf{Sub-CPMK-P8:} Melatih Random Forest; membandingkan metrik dengan baseline; membaca \texttt{feature\_importances\_} secara kritis (RPS Mg.\ 9).
\end{subcpmk}

\SesiRPS{9}{P8}{$\sim$170 menit}{CPMK-P3}

\section{Teori}

\begin{konsep}
Random Forest menggabungkan banyak pohon pada bootstrap sampel dan subset fitur acak; mengurangi varians dibanding satu pohon dalam. \texttt{feature\_importances\_} adalah proxy berbasis penurunan impurity---bukan kausalitas.
\end{konsep}

\section{Contoh kode}

\begin{lstlisting}[caption=RandomForestClassifier pada Iris]
from sklearn.ensemble import RandomForestClassifier
from sklearn.datasets import load_iris
from sklearn.model_selection import train_test_split

X, y = load_iris(return_X_y=True)
X_train, X_test, y_train, y_test = train_test_split(
    X, y, test_size=0.25, random_state=42, stratify=y
)
rf = RandomForestClassifier(n_estimators=200, random_state=42)
rf.fit(X_train, y_train)
print("RF score", rf.score(X_test, y_test))
\end{lstlisting}

\begin{lstlisting}[caption=Importansi fitur (urutkan)]
import pandas as pd
imp = pd.Series(rf.feature_importances_)
print(imp.sort_values(ascending=False))
\end{lstlisting}

\begin{lstlisting}[caption=RandomForestRegressor (diabetes)]
from sklearn.ensemble import RandomForestRegressor
from sklearn.datasets import load_diabetes
X, y = load_diabetes(return_X_y=True)
X_train, X_test, y_train, y_test = train_test_split(X, y, test_size=0.2, random_state=42)
rfr = RandomForestRegressor(n_estimators=200, random_state=42)
rfr.fit(X_train, y_train)
print("R2 test", rfr.score(X_test, y_test))
\end{lstlisting}

\begin{lstlisting}[caption=Catatan interpretasi di Markdown]
# Importansi tinggi != fitur penyebab; bisa karena korelasi atau leakage
\end{lstlisting}

\section{Referensi}

\begin{itemize}[leftmargin=*]
  \item sklearn ensemble. \url{https://scikit-learn.org/stable/modules/ensemble.html}
\end{itemize}

\begin{asesmen}
\textbf{Formatif:} perbandingan RF vs model minggu lalu; plot atau tabel importansi.
\end{asesmen}

\begin{rangkuman}
Minggu 9 memperkenalkan ensemble pohon untuk akurasi dan stabilitas relatif.

\textbf{Kata kunci:} RandomForest, feature importances
\end{rangkuman}

\end{document}
